\section{Literature Review}

One important stream in the literature uses recurrent neural networks and decomposition methods to improve forecasting quality. Fang et al\. propose a deep learning framework that combines improved singular spectrum analysis with a recurrent neural network to support adaptive decomposition and multi-feature fusion for energy consumption forecasting\cite{Fang2023}. Although their paper is about energy consumption rather than electricity price, it is relevant because it shows that recurrent models become much stronger when they are combined with preprocessing and feature-fusion techniques rather than used alone. Similarly, another paper combines EEMD, MRMR feature ranking, and BiLSTM, and reports improved forecasting performance in both 1-hour and 24-hour settings\cite{Shao2021}. Together, these studies suggest that recurrent architectures still matter, especially when they are paired with strong feature extraction.

This same logic appears in newer electricity-price work. Aksu et al\. propose a two-stage day-ahead EPF framework in which the signal is first decomposed with EMD and VMD, after which GRU is used for high-frequency components and LSTM for the remaining components\cite{Aksu2026}. Their hybrid model reportedly reduces RMSE by about 27.15\% compared with traditional machine-learning baselines and by 28.24\% compared with LSTM-based baselines. This is important because it suggests that decomposition is not just a preprocessing convenience, but a way to isolate rapid price fluctuations and improve prediction accuracy in volatile markets.

Older work also supports frequency-aware modeling. A 2008 paper propose an adaptive wavelet neural network for electricity-market price forecasting and report better performance than several competing neural and statistical baselines\cite{4494757}. A more recent paper reinforces preprocessing as an integral step\cite{Abedinia2017}. By combining neural networks with a metaheuristic optimization method, the authors were able to improve training and avoid poor local minima. These papers are not identical to the OpenRemote use case, but they reinforce a broader theme in the literature: wavelet-based decomposition and optimization-enhanced learning are useful when the underlying series is nonlinear, noisy, and unstable.

A second major direction is the use of attention and Transformer-based architectures. Wang et al\. present MultiDeT, a Transformer-based method for multi-energy load forecasting in integrated energy systems\cite{Wang2026}. Their model uses a one-encoder, multi-decoder structure so that different forecasting tasks can share one representation while learning different task-specific outputs. Even though OpenRemote does not necessarily need multiple decoders, the paper is still relevant because it shows how attention mechanisms can model interactions among multiple related variables. In other words, the core contribution is not only the multi-task architecture itself, but also the idea that one shared encoder can learn a useful multivariate representation.

A more directly relevant paper proposes a modified multi-head Transformer for multivariable building-energy forecasting\cite{Oliveira2023}. The authors compare the Transformer against LSTM and GRU baselines and report a lower MAPE for the Transformer-based model. This is useful for the current project because it supports the idea that attention-based architectures may outperform standard recurrent baselines when multiple input variables need to be combined effectively. The paper also makes an important practical point: LSTM and GRU remain meaningful baselines, but they may not be the best final choice when variable interaction becomes central to the task.

The most relevant recent paper for the OpenRemote use case is Temporal feature mixed inverted transformer: An inverted transformer for effective real-time electricity price forecasting\cite{Wang2026}. This study proposes a TFM-iTransformer model that combines an iTransformer backbone with volatility indicators, rate-of-change information, and future covariates. The paper is highly relevant because it explicitly addresses two problems that also apply to OpenRemote: electricity prices are highly volatile, and many forecasting models still underuse known future variables such as weather or dispatch-related information. The reported contribution is therefore not just architectural novelty, but a better way of modeling price uncertainty together with multivariate future information.

Long-horizon forecasting is another recurring concern. Ma and Zhang improve Autoformer by introducing a multi-scale feature fusion network and date-time encoding, producing the MD-Autoformer model\cite{Ma2025}. Their experiments show that the modified model improves on Autoformer across several datasets, but they also note that performance becomes difficult at prediction length 192. This is especially significant for the OpenRemote problem, because 48 hours at 15-minute intervals also equals 192 prediction steps. The implication is that long-range forecasting is possible, but it remains challenging even for advanced Transformer-like models.

Hybrid architectures are also promising. Al-Ali et al\. combine CNN, LSTM, and Transformer blocks for solar-energy production forecasting and report that the hybrid outperforms several simpler model combinations\cite{math11030676}. Another paper that employes ensemble models proposes a VMD-based multi-attention feature fusion model for electricity price forecasting that integrates GRU, TCN, SENet, and multi-head attention, and they report strong results against several comparison models\cite{Xu2025}. These studies suggest that the best-performing energy-forecasting models increasingly rely on architectural combinations that let different modules specialize in short-term fluctuations, periodic patterns, and feature weighting.

Finally, the broader survey literature supports these patterns. Kong et al\. argue that strong time-series forecasting models increasingly depend on external variables, multi-horizon design, and better feature extraction methods\cite{Kong2025}. Chen et al\. similarly propose a multivariate forecasting model that explicitly combines temporal features, periodicity, and cross-variable spatial relationships through LSTM, attention, graph-based periodic modeling, and CNNs\cite{CHEN2025126302}. Manero et al\., in a literature review on wind forecasting, also conclude that CNN and RNN/LSTM approaches outperform simple MLP approaches, especially when exogenous environmental variables are available\cite{Manero2018}. Across these sources, the same conclusion appears repeatedly: energy forecasting improves when models can exploit multiple correlated inputs rather than treating the target series in isolation. 

However, there is a considerable gap on the explicit inclusion of weather data for EPF\. A 2016 survey on Wind Energy, Load and Price forecasting mentions weather data as `[playing] a vital role in load forecasting', yet their approach has a higher degree of codependency between variables\cite{7754792}.
Weather patters are more linearly correlated to energy load than prices (see Fig.\ref{fig.corrHeat}). Sgarlato et al\. strongly show that the inclusion of weather predictions has strong explanatory power\cite{9788043}. This finding is key
for our mission of expanding the forecasting window to 48h. However, it is worth noting that this correlation is tied to location. Other papers discuss the use of weather data for energy forecasting\cite{7903735,Brecl2018-zt}. Despite these papers supporting the inclusion of weather data,
they use older methodology, or simply limit themselves to deterministic models.

Lastly, no paper included xLSTM\cite{xLSTM2024} in their methodology, as expected when such is a recent technology. Despite its immediate use being for NLP, it has been used sucesfully for time-series predictions\cite{xLSTMTimeSeries2025,xLSTMMixer}. Its use could limit computational power while conserving the benefits of attention
used in Transformers.